%=========== ΛΥΣΗ - ΘΕΜΑ Γ ===========

\begin{thema}{Γ}
\begin{erwthma}
\item Με βάση την περιγραφή της γραφικής παράστασης της $f''(x)$, η δεύτερη παράγωγος ορίζεται ως τμηματικά ορισμένη συνάρτηση ως εξής:
\[ f''(x) = \begin{cases} 1, & x < 2 \\ 0, & x \ge 2 \end{cases} \]

\item Εφόσον η $f'$ είναι αρχική της $f''$ και ισχύει $f'(0) = 0$, για $x < 2$ έχουμε:
\[ f'(x) = f'(0) + \int_{0}^{x} f''(t) \d t = 0 + \int_{0}^{x} 1 \d t = [t]_{0}^{x} = x \]
Δεδομένου ότι η $f$ είναι δύο φορές παραγωγίσιμη στο $\mathbb{R}$, η $f'$ πρέπει να είναι συνεχής στο $x=2$. Επομένως:
\[ f'(2) = \lim_{x \to 2^-} f'(x) = \lim_{x \to 2^-} x = 2 \]
Για $x \ge 2$, έχουμε:
\[ f'(x) = f'(2) + \int_{2}^{x} f''(t) \d t = 2 + \int_{2}^{x} 0 \d t = 2 + 0 = 2 \]
Συνεπώς, ο τύπος της $f'(x)$ είναι:
\[ f'(x) = \begin{cases} x, & x < 2 \\ 2, & x \ge 2 \end{cases} \]
Η συνάρτηση $f'(x)$ είναι $\textbf{γνησίως αύξουσα}$ στο $(-\infty, 2)$ και $\textbf{σταθερή}$ στο $[2, +\infty)$. Επειδή $\lim_{x \to -\infty} f'(x) = -\infty$, η $f'(x)$ δεν παρουσιάζει ελάχιστο στο $\mathbb{R}$.

\item Για $x < 2$, χρησιμοποιώντας το $f(0) = 0$:
\[ f(x) = f(0) + \int_{0}^{x} f'(t) \d t = 0 + \int_{0}^{x} t \d t = \left[ \frac{1}{2}t^2 \right]_{0}^{x} = \frac{1}{2}x^2 \]
Λόγω της παραγωγισιμότητας της $f$ στο $\mathbb{R}$, η $f$ είναι συνεχής στο $x=2$, άρα:
\[ f(2) = \lim_{x \to 2^-} f(x) = \lim_{x \to 2^-} \frac{1}{2}x^2 = \frac{1}{2}(2)^2 = 2 \]
Για $x \ge 2$, έχουμε:
\[ f(x) = f(2) + \int_{2}^{x} f'(t) \d t = 2 + \int_{2}^{x} 2 \d t = 2 + [2t]_{2}^{x} = 2 + (2x - 4) = 2x - 2 \]
Άρα, ο τύπος της $f(x)$ είναι:
\[ f(x) = \begin{cases} \frac{1}{2}x^2, & x < 2 \\ 2x - 2, & x \ge 2 \end{cases} \]
Για τη συνέχεια στο $x=2$:
\[ \lim_{x \to 2^-} f(x) = \frac{1}{2}(2)^2 = 2, \quad \lim_{x \to 2^+} f(x) = 2(2) - 2 = 2, \quad f(2) = 2 \]
Εφόσον $\lim_{x \to 2^-} f(x) = \lim_{x \to 2^+} f(x) = f(2)$, η συνάρτηση $f$ είναι $\textbf{συνεχής}$ στο $x=2$.

\item Μελετάμε το πρόσημο της $f'(x)$:
\begin{itemize}
    \item $f'(x) = 0 \Leftrightarrow x = 0$
    \item $f'(x) < 0 \Leftrightarrow x < 0$
    \item $f'(x) > 0 \Leftrightarrow x > 0$
\end{itemize}
Συνεπώς:
\begin{itemize}
    \item Στο διάστημα $(-\infty, 0)$, η $f'(x) < 0$, άρα η $f$ είναι $\textbf{γνησίως φθίνουσα}$.
    \item Στο διάστημα $(0, +\infty)$, η $f'(x) > 0$ (καθώς $f'(x)=x$ για $0 < x < 2$ και $f'(x)=2$ για $x \ge 2$), άρα η $f$ είναι $\textbf{γνησίως αύξουσα}}$.
\end{itemize}
Επειδή η $f'(x)$ αλλάζει πρόσημο από αρνητικό σε θετικό στο $x=0$, η συνάρτηση $f$ παρουσιάζει $\textbf{τοπικό ελάχιστο}$ στο $x=0$ με τιμή $f(0) = 0$.
\end{itemize}
\end{erwthma}
\end{thema}

%=========== ΤΕΛΟΣ ΛΥΣΗΣ - ΘΕΜΑ Γ ===========